%%%%%%%%%%%%%%%%%%%%%%%%%%%%%%%%%%%%%%%%%%%%%%%%%%%%%%%%%%%%%%%%%%%%%%%%%%%
% Chapter: Bridge Circuits
% Falstad examples: examples/bridges (generated by falstad/circuits/bridges.py)
% Capstone bridge (as in Capstone.tex): Pt100, 4.00 V, R1 = R3 = 3.9 kOhm,
% R4 = 100 Ohm, 0 .. 37.18 mV, gain 107.6
%%%%%%%%%%%%%%%%%%%%%%%%%%%%%%%%%%%%%%%%%%%%%%%%%%%%%%%%%%%%%%%%%%%%%%%%%%%

%%%%%%%%%%%%%%%%%%%%%%%%%%%%%%%%%%%%%%%%%%%%%%%%%%%%%%%%%%%%%%%%%%%%%%%%%%%
\section{The Wheatstone Bridge}
%%%%%%%%%%%%%%%%%%%%%%%%%%%%%%%%%%%%%%%%%%%%%%%%%%%%%%%%%%%%%%%%%%%%%%%%%%%
\nsl{Many sensors are resistors whose value changes a little with the
measured quantity: a platinum resistor changes by \SI{0.385}{\percent} per
kelvin, a strain gauge by \SI{0.2}{\percent} at the elastic limit of steel.
If we put such a sensor into a simple voltage divider, the interesting
change is a tiny ripple on a large constant voltage. Amplifying it would
amplify the constant part as well, and any drift of the supply would look
like a change of the measured quantity.
\newline
The bridge circuit solves this problem with a second voltage divider that
produces the constant part on its own. The difference between the two
divider outputs is zero when the sensor has its reference value and
contains only the change. The circuit was described by Christie in 1833
and made popular by Wheatstone in 1843 as a method to measure an unknown
resistance by adjusting a known one until the galvanometer between the two
dividers shows zero (\emph{null method}). Today the bridge is mostly used
in the \emph{deflection mode}: the resistors are fixed and the bridge
output voltage is amplified and digitised. This is the first stage of the
capstone task \cite{tietze2008,kester2005}.\newpage}

\begin{description}
\item[{\rot\bf Two voltage dividers}]~\\[-\lblskip]
\begin{center}
\begin{circuitikz}[scale=1.0, transform shape]
  \draw (0,0) to[V, invert, l=$V_B$] (0,4) -- (6,4);
  \draw (2.5,4) to[R, l_=$R_1$] (2.5,2) to[R, l_=$R_2$] (2.5,0);
  \draw (6,4) to[R, l=$R_3$] (6,2) to[R, l=$R_4$] (6,0) -- (0,0);
  \draw (2.5,4) node[circ]{} (2.5,0) node[circ]{};
  \draw (2.5,2) node[circ]{} node[left]{A} to[open, v^>=$V_d$] (6,2) node[circ]{} node[right]{B};
  \draw (0,0) node[ground]{};
\end{circuitikz}
\end{center}
\keyeq{Bridge output}{$\displaystyle V_d = V_A - V_B = V_B\left(\frac{R_2}{R_1+R_2}-\frac{R_4}{R_3+R_4}\right)$}
\nsl{Here $V_B$ is the excitation (supply) voltage of the bridge. Each
leg is an unloaded voltage divider from chapter~1; the output is the
difference of the two tap voltages. The output is taken between two nodes
of which neither is at ground: a \emph{differential} voltage.}
\os{\newpage}

\item[{\rot\bf Balance condition}]~\\[-\lblskip]
\bi
	\ite the bridge is \textbf{balanced} if $V_d = 0$:
\ei
\keyeq{Balance condition}{$\displaystyle \frac{R_1}{R_2} = \frac{R_3}{R_4} \quad\Leftrightarrow\quad R_1 R_4 = R_2 R_3$}
\bi
	\ite the balance does \textbf{not} depend on $V_B$: supply noise and drift do not appear at the output
	\ite null method: adjust a known resistor until $V_d = 0$, then $R_2 = R_1 R_4/R_3$
	\ite deflection method: fixed resistors, balanced at the reference point (e.g.\ \SI{0}{\celsius}), measure $V_d$
\ei
\nsl{At balance the output is zero, independent of the supply. Away from
balance the output is proportional to $V_B$: the bridge converts a
relative resistance change into a fraction of the excitation voltage.
This is why the excitation must be a stable reference, or the measurement
must be ratiometric (chapter~3).}
\os{\newpage}
\end{description}

%%%%%%%%%%%%%%%%%%%%%%%%%%%%%%%%%%%%%%%%%%%%%%%%%%%%%%%%%%%%%%%%%%%%%%%%%%%
\section{Output Voltage and Linearity}
%%%%%%%%%%%%%%%%%%%%%%%%%%%%%%%%%%%%%%%%%%%%%%%%%%%%%%%%%%%%%%%%%%%%%%%%%%%
\nsl{In most applications only one resistor of the bridge is a sensor
(\emph{quarter bridge}). We write its value as $R_2 = R_0(1+x)$, where
$x = \Delta R/R_0$ is the relative change, and choose the other resistors
such that the bridge is balanced at $x=0$. The output voltage is then a
function of $x$ only. For small $x$ it is proportional to $x$, for larger
$x$ it bends: the bridge is not linear. The reason is that the current in
the sensor leg changes with the sensor resistance.\newpage}

\begin{description}
\item[{\rot\bf Quarter bridge with equal resistors}]~\\[-\lblskip]
\bi
	\ite $R_1 = R_3 = R_4 = R_0$, sensor $R_2 = R_0(1+x)$:
	\[ V_d = V_B\left(\frac{1+x}{2+x}-\frac12\right) = \frac{V_B}{4}\cdot\frac{x}{1+x/2} \]
	\keyeq{Quarter bridge, small signal}{$\displaystyle V_d \approx \frac{V_B}{4}\,\frac{\Delta R}{R_0}$ \qquad relative error $\approx -\frac{x}{2}$}
	\ite example: $V_B = \SI{5}{\volt}$, $x = \SI{1}{\percent}$: $V_d = \SI{12.44}{\milli\volt}$ (estimate \SI{12.5}{\milli\volt})
	\ite $x = \SI{50}{\percent}$: $V_d = \SI{0.500}{\volt}$ instead of \SI{0.625}{\volt}
\ei
\falstad{wheatstone}
\os{\newpage}

\item[{\rot\bf The bridge is not linear}]~\\[-\lblskip]
\begin{center}
\begin{tikzpicture}
\begin{axis}[width=0.7\textwidth,height=55mm,xmin=-0.5,xmax=0.5,ymin=-0.17,ymax=0.17,
	xlabel={$x = \Delta R/R_0$},ylabel={$V_d/V_B$},axis lines=left,grid=major,
	font=\scriptsize,legend style={at={(0.03,0.97)},anchor=north west,draw=none},
	yticklabel style={/pgf/number format/fixed}]
	\addplot[rwucyandark,dashed,thick,domain=-0.5:0.5] {x/4};
	\addplot[rwuviolet,very thick,domain=-0.5:0.5,samples=50] {x/(2*(2+x))};
	\legend{linear: $x/4$, exact: $\frac{x}{2(2+x)}$}
\end{axis}
\end{tikzpicture}
\end{center}
\bi
	\ite positive $x$: output too small, negative $x$: too large
	\ite strain gauges ($x < \SI{0.5}{\percent}$): negligible; Pt100 over \SI{100}{\kelvin} ($x = \SI{38.5}{\percent}$): \SI{16}{\percent}
\ei
\os{\newpage}

\item[{\rot\bf Unequal arms: trade sensitivity for linearity}]~\\[-\lblskip]
\bi
	\ite $R_1 = R_3 = m\,R_0$, $R_4 = R_0$, sensor $R_0(1+x)$:
	\[ V_d = V_B\,\frac{m\,x}{(m+1)(m+1+x)} \approx V_B\,\frac{m}{(m+1)^2}\,x \]
	\ite sensitivity $m/(m+1)^2$: maximum $1/4$ at $m = 1$; $m = 39$: $0.024$
	\ite relative nonlinearity $\approx -x/(m+1)$: $m=39$ makes it 20 times smaller
	\ite large $m$ also limits the sensor current: $I = V_B/(R_1+R_\mathrm{sensor})$ (self-heating!)
\ei
\nsl{A large ratio $m$ turns the sensor leg into an almost ideal current
source: the current hardly changes with the sensor resistance, so the
voltage across the sensor is nearly proportional to its resistance. The
price is a smaller output. Since the amplifier needs a gain anyway, a few
tens of millivolts are easy to handle, while a nonlinearity would have to
be corrected in software. Other ways to linearise a bridge are a constant
current excitation or an op-amp that keeps the bridge at balance
\cite{kester2005}.}
\os{\newpage}
\end{description}

\exercise{Wheatstone bridge}
A bridge with $R_1 = R_3 = R_4 = \SI{1}{\kilo\ohm}$ is supplied with
$V_B = \SI{5}{\volt}$. The resistor $R_2 = R_x$ is the sensor.
\begin{talist}
\ta Calculate $V_d$ exactly and with the small-signal formula for $R_x = \SI{1010}{\ohm}$.
\ta Repeat for $R_x = \SI{1500}{\ohm}$. How large is the error of the formula?
\ta Check both values in Falstad.
\end{talist}

\begin{solution}
\begin{talist}
\ta $V_d = 5\,(1010/2010 - 0.5) = \SI{12.44}{\milli\volt}$; formula: $\frac{5}{4}\cdot 0.01 = \SI{12.5}{\milli\volt}$ (\SI{0.5}{\percent} too high, $\approx x/2$).
\ta $V_d = 5\,(1500/2500 - 0.5) = \SI{0.500}{\volt}$; formula: $\frac54\cdot 0.5 = \SI{0.625}{\volt}$: \SI{25}{\percent} too high.
\ta \falstad{wheatstone}: the voltmeter shows \SI{12.44}{\milli\volt}; with the slider at \SI{1500}{\ohm} it shows \SI{0.500}{\volt}.
\end{talist}
\end{solution}

%%%%%%%%%%%%%%%%%%%%%%%%%%%%%%%%%%%%%%%%%%%%%%%%%%%%%%%%%%%%%%%%%%%%%%%%%%%
\section{Half and Full Bridges: Strain Gauges}
%%%%%%%%%%%%%%%%%%%%%%%%%%%%%%%%%%%%%%%%%%%%%%%%%%%%%%%%%%%%%%%%%%%%%%%%%%%
\nsl{A strain gauge is a thin metal foil meander on a plastic carrier,
glued onto a mechanical part. When the part is stretched, the conductor
becomes longer and thinner and its resistance rises. The relative change
is proportional to the strain $\varepsilon = \Delta l/l$; the factor is the
\emph{gauge factor} $k$, about 2 for metal foil gauges. Typical values are
\SI{120}{\ohm}, \SI{350}{\ohm} and \SI{1}{\kilo\ohm}; a strain of
\SI{1000}{\micro\metre/\metre} (\SI{0.1}{\percent}, close to the elastic
limit of steel) changes the resistance by only \SI{0.2}{\percent}. On a
bending beam, gauges on the top are stretched and gauges on the bottom are
compressed by the same amount. Placing them in adjacent arms of the bridge
doubles or quadruples the signal and makes the bridge exactly linear. Load
cells of kitchen and industrial scales are full bridges.\newpage}

\begin{description}
\item[{\rot\bf Quarter, half and full bridge}]~\\[-\lblskip]
\begin{center}
\begin{circuitikz}[scale=0.72, transform shape]
  \def\brg#1#2#3#4#5#6{%
    \draw (#1,4) -- (#1+2.4,4);
    \draw (#1,4) to[R, l_=#2] (#1,2) to[R, l_=#3] (#1,0) -- (#1+2.4,0);
    \draw (#1+2.4,4) to[R, l=#4] (#1+2.4,2) to[R, l=#5] (#1+2.4,0);
    \draw (#1,2) node[circ]{} to[open, v^>=$V_d$] (#1+2.4,2) node[circ]{};
    \node[font=\Large] at (#1+1.2,-0.8) {#6};}
  \brg{0}{$R$}{$R(1{+}x)$}{$R$}{$R$}{quarter}
  \brg{6}{$R(1{-}x)$}{$R(1{+}x)$}{$R$}{$R$}{half}
  \brg{12}{$R(1{-}x)$}{$R(1{+}x)$}{$R(1{+}x)$}{$R(1{-}x)$}{full}
\end{circuitikz}
\end{center}
\bi
	\ite quarter: $V_d \approx \frac{V_B}{4}x$ \quad half: $V_d = \frac{V_B}{2}x$ \quad full: $V_d = V_B\,x$ \quad ($x = k\,\varepsilon$)
	\ite half and full bridge are exactly linear: the leg current stays constant
	\ite temperature changes all gauges equally: compensated
\ei
\falstad{bridge_types}
\nsl{\par In the half bridge the left leg contains $R(1-x)$ and $R(1+x)$; the
sum is $2R$, so the leg current does not change, and the tap voltage is
$V_B(1+x)/2$, exactly linear. In the full bridge the right leg moves in the
opposite direction and doubles the output once more. A temperature change
alters all four gauges by the same factor and leaves the ratios, and thus
the output, unchanged. In a quarter bridge the same effect is obtained with
an unstrained \emph{dummy gauge} on the same material as $R_4$.}
\os{\newpage}
\end{description}

\exercise{Strain gauge bridge}
Strain gauges with $R = \SI{350}{\ohm}$ and $k = 2$ are glued to a beam. The
strain is $\varepsilon = \SI{1000}{\micro\metre/\metre}$, the bridge supply
$V_B = \SI{5}{\volt}$.
\begin{talist}
\ta Calculate $\Delta R$.
\ta Calculate the output of a quarter, a half and a full bridge.
\ta Which gain is needed to get \SI{2}{\volt} at this strain from the full bridge?
\end{talist}

\begin{solution}
\begin{talist}
\ta $x = k\varepsilon = 2\cdot 10^{-3}$, $\Delta R = 0.002\cdot\SI{350}{\ohm} = \SI{0.7}{\ohm}$.
\ta Quarter: $V_d = 5\,(350.7/700.7 - 0.5) = \SI{2.4975}{\milli\volt}$ ($\approx\frac54\cdot 0.002 = \SI{2.5}{\milli\volt}$);
	half: $\frac52\cdot 0.002 = \SI{5.00}{\milli\volt}$; full: $5\cdot 0.002 = \SI{10.0}{\milli\volt}$.
\ta $G = \SI{2}{\volt}/\SI{10}{\milli\volt} = 200$, with an instrumentation amplifier (chapter~5).
\item[] \falstad{bridge_types}:\newline the three voltmeters show \SI{2.50}{\milli\volt}, \SI{5.00}{\milli\volt} and \SI{10.0}{\milli\volt}.
\end{talist}
\end{solution}

%%%%%%%%%%%%%%%%%%%%%%%%%%%%%%%%%%%%%%%%%%%%%%%%%%%%%%%%%%%%%%%%%%%%%%%%%%%
\section{Resistive Temperature Sensors}
%%%%%%%%%%%%%%%%%%%%%%%%%%%%%%%%%%%%%%%%%%%%%%%%%%%%%%%%%%%%%%%%%%%%%%%%%%%
\nsl{The resistance of metals rises with temperature. Platinum is used for
precision sensors because it is chemically stable and its characteristic
is reproducible and standardised (IEC~60751). A Pt100 has \SI{100}{\ohm}
at \SI{0}{\celsius}, a Pt1000 \SI{1000}{\ohm}. Between \SI{0}{\celsius} and
\SI{100}{\celsius} the mean temperature coefficient is
$\alpha = \SI{3.85e-3}{\per\kelvin}$, so a Pt100 changes by
\SI{0.385}{\ohm/\kelvin}. Over wider ranges the Callendar--Van Dusen
equation describes the slight curvature. Thermistors (NTC: negative
temperature coefficient) are semiconductor sensors: cheap, ten times more
sensitive, but strongly nonlinear and less accurate.\newpage}

\begin{description}
\item[{\rot\bf Platinum sensors Pt100 and Pt1000}]~\\[-\lblskip]
\bi
	\ite linear model ($\SI{0}{\celsius}\ldots\SI{100}{\celsius}$):
	\keyeq{Pt100 / Pt1000}{$\displaystyle R(T) = R_0\,(1+\alpha T)$, \quad $R_0 = \SI{100}{\ohm}$ / \SI{1000}{\ohm}, \quad $\alpha = \SI{3.85e-3}{\per\kelvin}$}
	\ite Pt100: \SI{100}{\ohm} at \SI{0}{\celsius}, \SI{138.5}{\ohm} at \SI{100}{\celsius}, \SI{0.385}{\ohm/\kelvin}
	\ite Callendar--Van Dusen ($T \geq \SI{0}{\celsius}$): $R(T) = R_0(1 + A T + B T^2)$
	\itee $A = \SI{3.9083e-3}{\per\kelvin}$, $B = \SI{-5.775e-7}{\per\kelvin\squared}$
	\ite tolerance class B: $\pm(\SI{0.3}{\kelvin} + 0.005\,|T|)$, class A: $\pm(\SI{0.15}{\kelvin} + 0.002\,|T|)$
\ei
\os{\newpage}

\item[{\rot\bf Pt100 characteristic}]~\\[-\lblskip]
\begin{center}
\begin{tikzpicture}
\begin{axis}[width=0.66\textwidth,height=50mm,xmin=-50,xmax=400,ymin=60,ymax=260,
	xlabel={$T$ / \si{\celsius}},ylabel={$R$ / \si{\ohm}},axis lines=left,grid=major,
	font=\scriptsize,legend style={at={(0.97,0.03)},anchor=south east,draw=none}]
	\addplot[rwucyandark,dashed,thick,domain=-50:400] {100*(1+0.00385*x)};
	\addplot[rwuviolet,thick,domain=0:400,samples=40] {100*(1+3.9083e-3*x-5.775e-7*x^2)};
	\addplot[rwuviolet,thick,domain=-50:0,samples=10] {100*(1+3.9083e-3*x-5.775e-7*x^2-4.183e-12*(x-100)*x^3)};
	\addplot[only marks,mark=*,mark size=1.5pt,rwuviolet] coordinates {(0,100) (100,138.51)};
	\legend{linear $\alpha = 0.00385$, Callendar--Van Dusen}
\end{axis}
\end{tikzpicture}
\end{center}
\bi
	\ite \SIrange{0}{100}{\celsius}: linear model and standard agree at both ends, \SI{0.4}{\kelvin} apart at \SI{50}{\celsius}
	\ite \SI{400}{\celsius}: \SI{254}{\ohm} (linear) vs.\ \SI{247}{\ohm}: \SI{18}{\kelvin} error
\ei
\nsl{For $T < \SI{0}{\celsius}$ the Callendar--Van Dusen equation has an
additional term $C\,(T-\SI{100}{\celsius})\,T^3$ with
$C = \SI{-4.183e-12}{\per\kelvin\tothe{4}}$. At \SI{50}{\celsius} the
standard gives \SI{119.40}{\ohm}, the linear model \SI{119.25}{\ohm}; the
difference of \SI{0.15}{\ohm} corresponds to \SI{0.4}{\kelvin}. For the
capstone with its resolution of \SI{12.5}{\kelvin} per LSB the linear model
is more than good enough.}
\os{\newpage}

\item[{\rot\bf NTC thermistors}]~\\[-\lblskip]
\begin{center}
\begin{tikzpicture}
\begin{axis}[width=0.55\textwidth,height=42mm,xmin=-20,xmax=100,ymin=0,ymax=70,
	xlabel={$T$ / \si{\celsius}},ylabel={$R$ / \si{\kilo\ohm}},axis lines=left,grid=major,font=\scriptsize]
	\addplot[rwuviolet,thick,domain=-20:100,samples=60] {10*exp(3950*(1/(x+273.15)-1/298.15))};
\end{axis}
\end{tikzpicture}
\end{center}
\bi
	\ite $R(T) = R_{25}\,\exp\!\left(B\left(\frac1T - \frac{1}{T_{25}}\right)\right)$, e.g.\ $R_{25} = \SI{10}{\kilo\ohm}$, $B = \SI{3950}{\kelvin}$ ($T$ in kelvin)
	\ite sensitivity at \SI{25}{\celsius}: $-B/T^2 = -\SI{4.4}{\percent/\kelvin}$ (Pt: $+\SI{0.39}{\percent/\kelvin}$)
	\ite strongly nonlinear: linearise with a series resistor, or a table in software
\ei
\os{\newpage}

\item[{\rot\bf Self-heating: why the sensor current is limited}]~\\[-\lblskip]
\bi
	\ite the measuring current heats the sensor: $P = I^2 R$
	\ite self-heating coefficient $E$: e.g.\ \SI{0.1}{\kelvin/\milli\watt} (thin film, in water) \ldots\ \SI{1}{\kelvin/\milli\watt} (in still air)
	\keyeq{Self-heating error}{$\Delta T = E\cdot I^2 R$}
	\ite Pt100, \SI{1}{\milli\ampere}: $P = \SI{0.1}{\milli\watt}$, $\Delta T = \SIrange{0.01}{0.1}{\kelvin}$
	\ite Pt100, \SI{10}{\milli\ampere}: $P = \SI{10}{\milli\watt}$, $\Delta T = \SIrange{1}{10}{\kelvin}$
	\ite rule: Pt100 $\leq \SI{1}{\milli\ampere}$, Pt1000 $\leq \SI{0.3}{\milli\ampere}$ \quad (capstone: $I \leq \SI{1}{\milli\ampere}$)
\ei
\nsl{The sensor measures its own temperature, not that of the medium. As
long as the self-heating is constant it could be calibrated out, but it
depends on how well the sensor is coupled to the medium: a sensor in
flowing water is cooled much better than the same sensor in still air. A
smaller current reduces the error with the square, but also the signal
linearly. For a Pt1000 the same power is reached at a $\sqrt{10}$ times
smaller current.}
\os{\newpage}
\end{description}

\exercise{Pt100 values}
\begin{talist}
\ta Calculate the resistance of a Pt100 at \SI{25}{\celsius} and \SI{37}{\celsius} with the linear model.
\ta Calculate $R(\SI{100}{\celsius})$ and $R(\SI{200}{\celsius})$ with the Callendar--Van Dusen equation. Which temperature does the linear model assign to these resistances?
\ta A Pt100 is read with \SI{5}{\milli\ampere} in still air ($E = \SI{0.5}{\kelvin/\milli\watt}$). How large is the self-heating error at \SI{0}{\celsius}?
\end{talist}

\begin{solution}
\begin{talist}
\ta $R(25) = 100\,(1+0.00385\cdot 25) = \SI{109.63}{\ohm}$, $R(37) = \SI{114.25}{\ohm}$.
\ta $R(100) = 100\,(1 + 0.39083 - 0.005775) = \SI{138.51}{\ohm}$,
	linear model: $T = 0.3851/0.00385 = \SI{100.0}{\celsius}$.
	$R(200) = 100\,(1+0.78166-0.0231) = \SI{175.86}{\ohm}$ $\Rightarrow$ linear: $T = \SI{197.0}{\celsius}$ (\SI{3}{\kelvin} too low).
\ta $P = (\SI{5}{\milli\ampere})^2\cdot\SI{100}{\ohm} = \SI{2.5}{\milli\watt}$, $\Delta T = \SI{1.25}{\kelvin}$: four times the class B tolerance at \SI{0}{\celsius} (\SI{0.3}{\kelvin}).
\end{talist}
\end{solution}

%%%%%%%%%%%%%%%%%%%%%%%%%%%%%%%%%%%%%%%%%%%%%%%%%%%%%%%%%%%%%%%%%%%%%%%%%%%
\section{Designing the Sensor Bridge}
%%%%%%%%%%%%%%%%%%%%%%%%%%%%%%%%%%%%%%%%%%%%%%%%%%%%%%%%%%%%%%%%%%%%%%%%%%%
\nsl{We now design the bridge of the capstone task. The requirements: a
Pt100 for \SIrange{0}{100}{\celsius}, a sensor current of at most
\SI{1}{\milli\ampere}, zero output at \SI{0}{\celsius}, and an amplifier
that produces exactly \SI{4.00}{\volt} at \SI{100}{\celsius}. The bridge is
excited from the \SI{4.00}{\volt} reference that also feeds the ladder of
the flash converter, so that the measurement is ratiometric (chapter~3).
The same design is used in the capstone chapter.\newpage}

\begin{description}
\item[{\rot\bf The capstone bridge}]~\\[-\lblskip]
\begin{center}
\begin{circuitikz}[scale=0.9, transform shape]
  \draw (0,0) to[V, invert, l=$\Vref{=}\SI{4.00}{\volt}$] (0,4) -- (6.5,4);
  \draw (3,4) to[R, l_=$R_1{=}\SI{3.9}{\kilo\ohm}$] (3,2) to[thR, l_=Pt100] (3,0);
  \draw (6.5,4) to[R, l=$R_3{=}\SI{3.9}{\kilo\ohm}$] (6.5,2) to[R, l=$R_4{=}\SI{100}{\ohm}$] (6.5,0) -- (0,0);
  \draw (3,4) node[circ]{} (3,0) node[circ]{};
  \draw (3,2) node[circ]{} node[below right]{$V_S$} to[open, v^>=$V_{BR}$] (6.5,2) node[circ]{} node[below left]{$V_R$};
  \draw (0,0) node[ground]{};
\end{circuitikz}
\end{center}
\bi
	\ite $R_4 = R_0 = \SI{100}{\ohm}$: balanced at \SI{0}{\celsius}
	\ite $I = \Vref/(R_1 + R) \leq \SI{1}{\milli\ampere}$ $\Rightarrow$ $R_1 \geq \SI{3.9}{\kilo\ohm}$; $R_3 = R_1$ ($m = 39$)
	\ite $V_{BR}(T) = \Vref\left(\frac{R(T)}{R_1+R(T)} - \frac{R_4}{R_3+R_4}\right)$
\ei
\os{\newpage}

\item[{\rot\bf Numbers of the capstone bridge}]~\\[-\lblskip]
\begin{center}
\begin{tabular}{lccc}
\toprule
 & \SI{0}{\celsius} & \SI{50}{\celsius} & \SI{100}{\celsius} \\
\midrule
sensor $R$ & \SI{100}{\ohm} & \SI{119.25}{\ohm} & \SI{138.5}{\ohm} \\
sensor current & \SI{1.000}{\milli\ampere} & \SI{0.995}{\milli\ampere} & \SI{0.990}{\milli\ampere} \\
$V_S$ & \SI{100.0}{\milli\volt} & \SI{118.7}{\milli\volt} & \SI{137.2}{\milli\volt} \\
$V_R$ & \SI{100.0}{\milli\volt} & \SI{100.0}{\milli\volt} & \SI{100.0}{\milli\volt} \\
$V_{BR}$ & \SI{0}{\milli\volt} & \SI{18.68}{\milli\volt} & \SI{37.18}{\milli\volt} \\
\bottomrule
\end{tabular}
\end{center}
\bi
	\ite sensitivity $\approx \SI{0.37}{\milli\volt/\kelvin}$; nonlinearity at \SI{50}{\celsius}: $+\SI{0.09}{\milli\volt}$ (\SI{0.25}{\percent} of full scale)
	\ite required gain: $G = \SI{4.00}{\volt}/\SI{37.18}{\milli\volt} = 107.6$
	\ite self-heating: $P = \SI{0.1}{\milli\watt}$
\ei
\falstad{pt100_bridge}
\os{\newpage}

\item[{\rot\bf Bridge output over temperature}]~\\[-\lblskip]
\begin{center}
\begin{tikzpicture}
\begin{axis}[width=0.66\textwidth,height=48mm,xmin=0,xmax=100,ymin=0,ymax=40,
	xlabel={$T$ / \si{\celsius}},ylabel={$V_{BR}$ / mV},axis lines=left,grid=major,font=\scriptsize,
	legend style={at={(0.03,0.97)},anchor=north west,draw=none}]
	\addplot[rwucyandark,dashed,thick,domain=0:100] {0.3718*x};
	\addplot[rwuviolet,thick,domain=0:100,samples=30] {4000*(100*(1+0.00385*x)/(3900+100*(1+0.00385*x)) - 0.025)};
	\legend{straight line through the end points, bridge ($m = 39$)}
\end{axis}
\end{tikzpicture}
\end{center}
\bi
	\ite with $m = 39$ the curve is almost straight: maximum deviation \SI{0.09}{\milli\volt} $\approx \SI{0.24}{\kelvin}$
	\ite with $m = 1$ (all arms \SI{100}{\ohm}) the deviation would be \SI{4.4}{\kelvin}
\ei
\os{\newpage}

\item[{\rot\bf Design card}]~\\[-\lblskip]
\designcard{Sensor bridge}{
\bi
	\ite excitation $V_B$: the ADC reference (ratiometric), here \SI{4.00}{\volt}
	\ite reference arm: $R_4 = R_\mathrm{sensor}(T_\mathrm{min})$ $\Rightarrow$ $V_d = 0$ at $T_\mathrm{min}$
	\ite current limit: $R_1 \geq V_B/I_\mathrm{max} - R_\mathrm{sensor}$, round up to E-series; $R_3 = R_1$
	\ite full-scale output $V_d(T_\mathrm{max})$ with the exact formula, gain $G = V_\mathrm{FS}/V_d(T_\mathrm{max})$
	\ite check nonlinearity ($\approx x/(m+1)$), self-heating ($I^2R\cdot E$), common-mode voltage ($\approx V_B/(m+1)$)
\ei}
\os{\newpage}
\end{description}

\exercise{Design of the capstone bridge}
A Pt100 ($\alpha = \SI{3.85e-3}{\per\kelvin}$) shall measure
\SIrange{0}{100}{\celsius}. The bridge is fed from $\Vref = \SI{4.00}{\volt}$,
the sensor current must not exceed \SI{1}{\milli\ampere}.
\begin{talist}
\ta Choose $R_1$, $R_3$ and $R_4$.
\ta Calculate the bridge output at \SI{0}{\celsius}, \SI{50}{\celsius} and \SI{100}{\celsius} and the sensor currents.
\ta How large is the deviation from a straight line at \SI{50}{\celsius}, in millivolt and in kelvin?
\ta Which total amplifier gain gives \SI{4.00}{\volt} at \SI{100}{\celsius}?
\end{talist}

\begin{solution}
\begin{talist}
\ta $R_4 = \SI{100}{\ohm}$ (balance at \SI{0}{\celsius}); $R_1 \geq \SI{4}{\volt}/\SI{1}{\milli\ampere} - \SI{100}{\ohm} = \SI{3.9}{\kilo\ohm}$ (E12), $R_3 = R_1 = \SI{3.9}{\kilo\ohm}$.
\ta $V_R = 4\cdot 100/4000 = \SI{100.0}{\milli\volt}$.
	\SI{0}{\celsius}: $V_{BR} = 0$, $I = \SI{1.000}{\milli\ampere}$.
	\SI{50}{\celsius}: $V_S = 4\cdot 119.25/4019.25 = \SI{118.68}{\milli\volt}$, $V_{BR} = \SI{18.68}{\milli\volt}$.
	\SI{100}{\celsius}: $V_S = 4\cdot 138.5/4038.5 = \SI{137.18}{\milli\volt}$, $V_{BR} = \SI{37.18}{\milli\volt}$, $I = \SI{0.990}{\milli\ampere}$.
\ta Straight line: $37.18/2 = \SI{18.59}{\milli\volt}$; deviation \SI{0.09}{\milli\volt}, with \SI{0.372}{\milli\volt/\kelvin}: \SI{0.24}{\kelvin}.
\ta $G = \SI{4.00}{\volt}/\SI{37.18}{\milli\volt} = 107.6$ (capstone: $10\times 10.76$).
\item[] \falstad{pt100_bridge}: with the slider at \SI{100}{\ohm} the voltmeter shows \SI{0}{\volt} and $V_R = \SI{100}{\milli\volt}$;
	at \SI{119.25}{\ohm} \SI{18.68}{\milli\volt}; at \SI{138.5}{\ohm} \SI{37.18}{\milli\volt}.
\end{talist}
\end{solution}

\exercise{Alternative: Pt1000 at 5 V}
A Pt1000 is used in a bridge with $V_B = \SI{5}{\volt}$, $R_1 = R_3 = \SI{4.7}{\kilo\ohm}$, $R_4 = \SI{1}{\kilo\ohm}$.
\begin{talist}
\ta Calculate the sensor current and the self-heating power at \SI{0}{\celsius} and \SI{100}{\celsius}.
\ta Calculate the bridge output at \SI{100}{\celsius} and the required gain for \SI{4.00}{\volt}.
\ta How large is the deviation from the straight line at \SI{50}{\celsius}? Compare with the Pt100 design.
\end{talist}

\begin{solution}
\begin{talist}
\ta $I = 5/5700 = \SI{0.877}{\milli\ampere}$, $P = I^2\cdot\SI{1}{\kilo\ohm} = \SI{0.77}{\milli\watt}$;
	at \SI{100}{\celsius}: $I = 5/6085 = \SI{0.822}{\milli\ampere}$, $P = \SI{0.94}{\milli\watt}$, almost ten times the Pt100.
\ta $V_d = 5\,(1385/6085 - 1000/5700) = \SI{260.9}{\milli\volt}$, $G = 4.00/0.2609 = 15.33$.
\ta $V_d(\SI{50}{\celsius}) = 5\,(1192.5/5892.5 - 0.17544) = \SI{134.7}{\milli\volt}$, straight line \SI{130.4}{\milli\volt}:
	\SI{4.3}{\milli\volt} $= \SI{1.6}{\percent}$ of full scale ($m = 4.7$), almost seven times worse than the Pt100 design with $m = 39$; in exchange the signal is seven times larger.
\item[] \falstad{pt1000_bridge}: voltmeter \SI{0}{\volt}, \SI{134.7}{\milli\volt} and \SI{260.9}{\milli\volt} at the slider positions 1000, 1192.5 and \SI{1385}{\ohm}.
\end{talist}
\end{solution}

%%%%%%%%%%%%%%%%%%%%%%%%%%%%%%%%%%%%%%%%%%%%%%%%%%%%%%%%%%%%%%%%%%%%%%%%%%%
\section{Reading the Bridge: Common-Mode Voltage}
%%%%%%%%%%%%%%%%%%%%%%%%%%%%%%%%%%%%%%%%%%%%%%%%%%%%%%%%%%%%%%%%%%%%%%%%%%%
\nsl{The bridge output is the difference of two node voltages, each of
which is far from zero. In the capstone bridge both nodes sit at about
\SI{100}{\milli\volt}; with $m = 1$ they would be at half the supply. The
mean of the two node voltages is the \emph{common-mode voltage}
$V_\mathrm{CM} = (V_S+V_R)/2$, the difference is the \emph{differential
voltage} $V_d = V_S - V_R$. The amplifier must amplify $V_d$ by about 100
and ignore $V_\mathrm{CM}$. An ordinary amplifier referenced to ground
cannot do this.\newpage}

\begin{description}
\item[{\rot\bf A single-ended amplifier fails}]~\\[-\lblskip]
\begin{center}
\begin{circuitikz}[scale=0.8, transform shape]
  \draw (0,0) to[V, invert, l=\SI{4}{\volt}] (0,4) -- (4,4);
  \draw (2,4) to[R, l_=\SI{3.9}{\kilo\ohm}] (2,2) to[thR, l_=Pt100] (2,0);
  \draw (4,4) to[R, l=\SI{3.9}{\kilo\ohm}] (4,2) to[R, l=\SI{100}{\ohm}] (4,0) -- (0,0);
  \draw (0,0) node[ground]{};
  \draw (4,2) node[circ]{} node[right]{$V_R$};
  \draw (2,2) node[circ]{} to[short,-o] (3,2) node[above]{$V_S$};
  \node[op amp, noinv input up, anchor=+] (A) at (7.6,2.7) {};
  \draw (6.6,2.7) node[ocirc]{} node[above]{$V_S$} -- (A.+);
  \draw (A.-) -- ++(-0.3,0) -- ++(0,-1.0) coordinate(m);
  \draw (m) to[R, l_=\SI{10}{\kilo\ohm}] ++(0,-1.9) node[ground]{};
  \draw (m) node[circ]{} to[R, l_=\SI{1066}{\kilo\ohm}] ++(3.3,0) coordinate(f) -- (f |- A.out);
  \draw (A.out) -- (f |- A.out) node[circ]{} to[short,-o] ++(0.8,0) node[right]{$\Vout$};
\end{circuitikz}
\end{center}
\bi
	\ite gain 107.6 applied to $V_S = \SI{100}{\milli\volt}$: $\Vout = \SI{10.75}{\volt}$ at \SI{0}{\celsius} (wanted: \SI{0}{\volt})
	\ite at \SI{100}{\celsius}: \SI{14.76}{\volt}; the useful signal is a small change on a large offset
	\ite any drift of $\Vref$ appears amplified at the output
\ei
\falstad{bridge_single_ended}
\os{\newpage}

\item[{\rot\bf Difference and instrumentation amplifier}]~\\[-\lblskip]
\begin{center}
\begin{circuitikz}[scale=0.7, transform shape]
  \node[op amp, noinv input up] (b1) at (1.2,3) {};
  \draw (b1.+) to[short,-o] ++(-0.6,0) node[left]{$V_R$};
  \draw (b1.-) -- ++(0,-0.6) -| (b1.out);
  \node[op amp, noinv input up] (b2) at (1.2,-0.6) {};
  \draw (b2.+) to[short,-o] ++(-0.6,0) node[left]{$V_S$};
  \draw (b2.-) -- ++(0,-0.6) -| (b2.out);
  \node[op amp] (da) at (7,1.2) {};
  \draw (b1.out) -- ++(0.4,0) |- ($(da.-)+(-2.4,0)$) to[R, l=\SI{10}{\kilo\ohm}] (da.-);
  \draw (b2.out) -- ++(0.4,0) |- ($(da.+)+(-2.4,0)$) to[R, l_=\SI{10}{\kilo\ohm}] (da.+);
  \draw ($(da.-)+(-0.4,0)$) node[circ]{} -- ++(0,1.2) coordinate(t)
        to[R, l=\SI{100}{\kilo\ohm}] (t -| da.out) -- (da.out);
  \draw ($(da.+)+(-0.4,0)$) node[circ]{} to[R, l=\SI{100}{\kilo\ohm}] ++(0,-1.8) node[ground]{};
  \draw (da.out) to[short,-o] ++(0.8,0) node[right]{$V_\mathrm{DIFF} = 10\,(V_S - V_R)$};
\end{circuitikz}
\end{center}
\bi
	\ite difference amplifier: amplifies $V_S - V_R$, rejects $V_\mathrm{CM}$ (if the resistors match)
	\ite buffers: no current from the bridge, the bridge is not loaded
	\ite second stage $G_2 = 10.76$: \SI{0}{\volt} \ldots\ \SI{4.00}{\volt} \quad (details: chapter~5)
\ei
\falstad{bridge_ina}
\nsl{\par Without the buffers the difference amplifier would load the bridge
with its input resistors. Here the bridge has a source resistance of only
$\SI{3.9}{\kilo\ohm}\parallelto\SI{100}{\ohm} = \SI{97.5}{\ohm}$, so
\SI{10}{\kilo\ohm} inputs would cause about \SI{1}{\percent} gain error,
and worse, a mismatch between the two inputs that converts common-mode
voltage into an output error. The two buffers and the difference amplifier
together form the classic three-op-amp instrumentation amplifier (INA),
treated in chapter~5. In Falstad the complete front end delivers
\SI{0}{\volt}, \SI{2.01}{\volt} and \SI{4.00}{\volt} at \SI{0}{\celsius},
\SI{50}{\celsius} and \SI{100}{\celsius}.}
\os{\newpage}
\end{description}

%%%%%%%%%%%%%%%%%%%%%%%%%%%%%%%%%%%%%%%%%%%%%%%%%%%%%%%%%%%%%%%%%%%%%%%%%%%
\section{Interference and Lead Resistance}
%%%%%%%%%%%%%%%%%%%%%%%%%%%%%%%%%%%%%%%%%%%%%%%%%%%%%%%%%%%%%%%%%%%%%%%%%%%
\nsl{The sensor is often far away from the electronics, connected by a
cable of several metres. The cable brings two problems. First, the mains
wiring and the fluorescent lamps in the lab create \SI{50}{\hertz}
electric and magnetic fields. The loop formed by the two sensor wires acts
as the secondary winding of a transformer, and the induced voltage appears
in series with the sensor: a differential signal that the instrumentation
amplifier amplifies just like the measurement signal. Interference that
couples equally into both bridge nodes is common-mode and rejected. In
the capstone the overhead lamps are modelled by a \SI{5}{\milli\volt},
\SI{50}{\hertz} source in series with the sensor. Second, the resistance of
the cable adds to the sensor resistance. For a Pt100 one ohm corresponds to
\SI{2.6}{\kelvin}.\newpage}

\begin{description}
\item[{\rot\bf 50\,Hz in series with the sensor}]~\\[-\lblskip]
\begin{center}
\begin{circuitikz}[scale=0.85, transform shape]
  \draw (0,0) to[V, invert, l=\SI{4}{\volt}] (0,4.5) -- (6,4.5);
  \draw (2.5,4.5) to[R, l_=\SI{3.9}{\kilo\ohm}] (2.5,2.7) to[thR, l_=Pt100] (2.5,1.3)
        to[sV, l_=\SI{5}{\milli\volt}\,\SI{50}{\hertz}] (2.5,0);
  \draw (6,4.5) to[R, l=\SI{3.9}{\kilo\ohm}] (6,2.7) to[R, l=\SI{100}{\ohm}] (6,0) -- (0,0);
  \draw (0,0) node[ground]{};
  \draw (2.5,2.7) node[circ]{} node[right]{$V_S$} (6,2.7) node[circ]{} node[left]{$V_R$};
\end{circuitikz}
\end{center}
\bi
	\ite at $V_S$: $\SI{5}{\milli\volt}\cdot\frac{R_1}{R_1+R} = \SI{5}{\milli\volt}\cdot\frac{3900}{4000} = \SI{4.88}{\milli\volt}$ amplitude
	\ite after $G = 107.6$: $\SI{0.52}{\volt}$ amplitude: about one LSB (\SI{0.5}{\volt})
	\ite remedies: twisted pair, shield, short leads, low-pass filter before the ADC (chapters~6, 7)
\ei
\falstad{bridge_50hz}
\nsl{\par Twisting the two sensor wires reduces the loop area and makes the
induced voltages of neighbouring twists cancel. A shield connected to
ground at the amplifier side diverts capacitively coupled currents. What
remains must be removed by filtering: the temperature changes slowly, so
everything above a few hertz can be cut off. Chapter~6 designs the
anti-aliasing filter and chapter~7 shows what happens if the interference
reaches the sampler unfiltered.}
\os{\newpage}

\item[{\rot\bf Lead resistance: 2-wire and 3-wire connection}]~\\[-\lblskip]
\begin{center}
\begin{circuitikz}[scale=0.72, transform shape]
  % 2-wire
  \draw (0,5) -- (3.5,5);
  \draw (0,5) to[R, l_=$R_1$] (0,3.2) to[R, l_=$R_L$] (0,1.8) to[thR, l_=Pt100] (0,0.4)
        to[R, l_=$R_L$] (0,-1) -- (3.5,-1);
  \draw (3.5,5) to[R, l=$R_3$] (3.5,3.2) to[R, l=$R_4$] (3.5,-1);
  \draw (0,3.2) node[circ]{} to[open, v^>=$V_d$] (3.5,3.2) node[circ]{};
  \node[font=\large] at (1.75,-1.7) {2-wire};
  % 3-wire
  \draw (8,5) -- (11.5,5);
  \draw (8,5) to[R, l_=$R_1$] (8,3.2) to[R, l_=$R_{L1}$] (8,1.8) to[thR, l_=Pt100] (8,0.4) -- (9.75,0.4);
  \draw (11.5,5) to[R, l=$R_3$] (11.5,3.2) to[R, l=$R_4$] (11.5,1.6) to[R, l=$R_{L2}$] (11.5,0.4) -- (9.75,0.4);
  \draw (9.75,0.4) node[circ]{} to[R, l=$R_{L3}$] (9.75,-1) node[ground]{};
  \draw (8,3.2) node[circ]{} to[open, v^>=$V_d$] (11.5,3.2) node[circ]{};
  \node[font=\large] at (9.75,-1.9) {3-wire};
\end{circuitikz}
\end{center}
\bi
	\ite 2-wire: $2R_L$ in series with the sensor; $R_L = \SI{0.5}{\ohm}$: \SI{1}{\ohm} $\hat{=}$ \SI{2.6}{\kelvin}
	\ite 3-wire: $R_{L1}$ in the sensor arm, $R_{L2}$ in the neighbouring arm $R_4$: they cancel at balance;
	$R_{L3}$ only carries the supply current
	\ite 4-wire: current source through two wires, voltage sensed by the other two (precision instruments)
\ei
\falstad{bridge_3wire}
\os{\newpage}
\end{description}

\exercise{Interference and lead resistance}
The capstone bridge ($\Vref = \SI{4}{\volt}$, $R_1=R_3=\SI{3.9}{\kilo\ohm}$,
$R_4 = \SI{100}{\ohm}$, Pt100) is followed by an amplifier with gain 107.6.
\begin{talist}
\ta A \SI{5}{\milli\volt}, \SI{50}{\hertz} interference voltage is induced in series with the sensor.
	Calculate its amplitude at $V_S$ and at the amplifier output. Why does the instrumentation amplifier not reject it?
\ta The sensor is connected with two leads of \SI{0.5}{\ohm} each. Which bridge output and temperature error result at \SI{0}{\celsius}?
\ta Repeat b) for the 3-wire connection, at \SI{0}{\celsius} and \SI{100}{\celsius}.
\end{talist}

\begin{solution}
\begin{talist}
\ta The source sees $R_1$ above and the sensor below: $\SI{5}{\milli\volt}\cdot 3900/4000 = \SI{4.875}{\milli\volt}$ at $V_S$,
	$\SI{4.875}{\milli\volt}\cdot 107.6 = \SI{0.52}{\volt}$ at the output. It changes only $V_S$, not $V_R$: it is a differential signal.
\ta $R = \SI{101}{\ohm}$: $V_d = 4\,(101/4001 - 0.025) = \SI{0.975}{\milli\volt}$; with \SI{0.372}{\milli\volt/\kelvin}: \SI{2.6}{\kelvin}.
\ta \SI{0}{\celsius}: both lower arms have \SI{100.5}{\ohm}, $V_d = 0$. \SI{100}{\celsius}: $V_d = \SI{37.16}{\milli\volt}$
	instead of \SI{37.18}{\milli\volt} (2-wire: \SI{38.13}{\milli\volt}); the remaining error is \SI{0.05}{\kelvin}.
\item[] \falstad{bridge_50hz}: the scope shows \SI{4.88}{\milli\volt} amplitude at $V_S$ and \SI{0.52}{\volt} at \texttt{Vout};
	\falstad{bridge_3wire}: voltmeters \SI{0.975}{\milli\volt} (2-wire) and \SI{0}{\volt} (3-wire) at \SI{0}{\celsius},
	\SI{38.13}{\milli\volt} and \SI{37.16}{\milli\volt} with both sliders at \SI{138.5}{\ohm}.
\end{talist}
\end{solution}

%%% Local Variables:
%%% mode: latex
%%% TeX-master: "docu"
%%% End:
